\documentclass[11pt, twocolumn, landscape, a4paper]{article}

\usepackage[margin=1cm]{geometry}
\usepackage[utf8]{inputenc}
\usepackage[spanish]{babel}
\usepackage{amssymb, amsmath, amsbsy}
\usepackage{tasks}
\usepackage{enumerate}
\usepackage{graphicx}

\usepackage[pdftex]{hyperref}
\hypersetup{pdftitle={Repaso en Aritmética},pdfauthor={Luis Carrillo Gutiérrez}}

\fontfamily{cmss}\selectfont
\renewcommand{\familydefault}{\sfdefault}
\pagestyle{empty}

\begin{document}
\subsubsection*{Cálculo de los divisores de un número natural ($\mathbb{N}$)}

\noindent Vamos a desarrollar un método para calcular los divisores de un número, por ejemplo de $72$. Primero se descompone el número $72$ en factores primos.

\begin{center}
\begin{tabular}{l|r}
72 & 2 \\
36 & 2 \\
18 & 2 \\
9 & 3 \\
3 & 3 \\
1 
\end{tabular}
\end{center}

$$72 = 2^{\,3} \cdot 3^{\,2}$$

Para conocer el número de divisores de $72$, se multiplican los {\tt exponentes} de los factores primos de su descomposición factorial, \emph{aumentados en una unidad}. 

$$
(3 + 1) \cdot (2 + 1) = 4 \cdot 3 = 12
$$

El número $72$ tendrá $12$ divisores.

% A continuación se verá cuáles son esos divisores.
\end{document}